% !TeX root = ../深度学习与大语言模型技术原理.tex
% 由主文档通过 \input 载入；请编译主文档。
\setcounter{chapter}{12} % 保留原书章号。
\chapter{拒绝筛选、蒸馏与推理时计算}
\label{ch:distillation}

\section{采样温度、核采样与拒绝筛选}

温度 $T>0$ 把 logits $z_j$ 变为 $p_T(j)\propto\exp(z_j/T)$。降低温度通常使分布集中，$T\to0$ 时质量趋向最大得分集合；遇到并列最大值时仍需规定选择规则。Top-$k$ 采样保留概率最大的 $k$ 个词元；核采样（nucleus sampling, top-$p$ sampling）保留累计概率至少为阈值 $p$ 的最小排序前缀，再重新归一化。这些操作改变生成分布，不能在使用过滤后采样结果时仍以未过滤分布作精确概率计算。

对提示生成多个回答，再用验证器保留高质量结果，是训练数据中的拒绝筛选。若原生成分布为 $q(y\mid x)$、通过概率为 $a(x,y)\in[0,1]$，且通过率非零，保留样本服从
\begin{equation}
q_{\mathrm{keep}}(y\mid x)=\frac{q(y\mid x)a(x,y)}{\sum_{y'}q(y'\mid x)a(x,y')}.
\label{eq:filter-distribution}
\end{equation}
这说明筛选偏向容易通过规则的回答。只有按特定接受率构造时，它才等价于经典概率论中从指定目标分布精确抽样的拒绝采样（rejection sampling）。行业中对候选筛选也使用同名术语，应明确语境。

若每题保留一个最佳结果，形成 Best-of-$N$ 数据，分布还依赖候选间的比较。增加 $N$ 既提高发现正确答案的机会，也提高利用评分器漏洞的机会。必须同时记录生成成本、通过率、题目覆盖和筛选后长度分布。

\section{软目标蒸馏与行为蒸馏}

\emph{知识蒸馏（knowledge distillation）}让学生模型学习教师行为。设教师、学生在同一上下文和词表上的温度分布为 $p_T,q_T$，一个常见目标为
\[
\mathcal L=(1-\lambda)\mathcal L_{\mathrm{hard}}
+\lambda T^2D_{\mathrm{KL}}(p_T\|q_T),\qquad 0\le\lambda\le1.
\]
软分布同时提供非目标词元的相对信息。$T^2$ 常用于补偿温度对梯度尺度的影响，但其效果依赖得分尺度与温度区间；并不是保证所有温度下梯度完全相同的恒等式\cite{distill}。

响应蒸馏只训练教师生成的离散文本，因而相当于在教师样本上做 SFT。推理轨迹蒸馏进一步保留中间步骤；工具行为蒸馏保留动作、参数和观察。教师写出的解释是可学习的输出轨迹，不应被等同于其内部计算过程。不同词表模型之间不能直接逐编号比较 logits，通常需要文本级训练或显式对齐。

离策略蒸馏（off-policy distillation）使用预先生成的教师轨迹；同策略蒸馏（on-policy distillation）让学生先产生自己常访问的前缀，再由教师在这些前缀上提供信号。后者可以减轻前缀分布失配，但增加在线教师调用与训练调度成本。学生可能缺少教师的表示容量，故即使目标一致也未必达到教师性能。

\section{提示设计、少样本示例与推理轨迹}

提示工程（prompt engineering）在固定模型参数下设计任务说明、输入边界、示例和输出约束。零样本（zero-shot）提示不给已解示例，少样本（few-shot）提示提供少量输入--输出对，帮助模型识别任务和格式。这种上下文学习（in-context learning）通过条件输入改变输出，不等于执行梯度更新；示例的顺序、标签和长度都可能影响结果，应在独立任务上评估。

思维链（chain of thought, CoT）在输出中加入中间推导或分步说明，使模型可以利用已生成的中间文本继续计算。它可能帮助需要多步组合的任务，也可能引入冗长错误和错误前提传播。可见推理轨迹是输出内容，不是模型内部计算机制的完整记录；其正确性需要数学核查、执行反馈或其他外部证据。

提示设计应优先明确问题、条件和可核验输出，再决定是否需要中间步骤。数据中的长推理轨迹可以用于蒸馏，但应检查每步是否有效、最终答案是否被泄漏以及是否包含与任务无关的固定格式。推理长度是资源消耗指标之一，不能单独作为质量标签。

\section{自一致性与验证器选择}

\emph{推理时计算扩展（test-time scaling）}在给定模型参数下，通过更多采样、搜索或验证改善结果。自一致性（self-consistency）从多条推理路径提取标准化答案，再投票\cite{selfconsistency}；Best-of-$N$ 用验证器选择高分候选。前者依赖错误是否分散，后者依赖验证器能否区分正确与错误。若模型系统性偏向同一错误，多数投票会加强错误。

作为理想参照，若各次生成独立且成功概率为 $p$，至少一次成功的概率为 $1-(1-p)^N$。但把成功候选识别出来还需额外选择机制，所以该概率通常是带完美识别器的可达上界，而不是实际交付答案的准确率。固定总计算预算时，更多短候选与更少长候选应分别比较。

思考预算（thinking budget）可以按任务难度、当前不确定性和剩余资源动态分配。一个决策目标是最大化 $\mathbb E[\text{正确性}]-\lambda\mathbb E[\text{成本}]$，成本可以是词元、时间或费用。困难度估计、停止规则和预算调度都需独立验证，不能把更长推理文本当成更强推理的直接证据。

\section{束搜索、树搜索与停止条件}

束搜索（beam search）保留得分较高的若干前缀，扩展后再截断。它近似优化序列概率或经长度调整的分数，并不直接优化事实正确性；如果概率模型偏好错误模式，扩大束宽也不能修复目标。验证器引导搜索则把任务反馈引入节点评分，但需要承担节点生成和评分成本。

蒙特卡罗树搜索（Monte Carlo tree search, MCTS）风格方法在选择、扩展、模拟和回传之间循环。例如 UCT 选择准则可写为 $\widehat Q(s,a)+c\sqrt{\log N(s)/N(s,a)}$，未访问动作另外优先探索。语言推理中，一个动作可能是一词元、一步推导或一次工具调用；动作粒度会显著改变树宽、深度和价值估计难度。采用相似公式并不自动获得经典有限状态搜索条件下的保证。

搜索应有明确的总预算、重复状态检测和终止验证。若验证器只检验最终答案，错误中间前缀可能浪费大量扩展；若用过程模型剪枝，误判又可能提前删除正确路径。实验应报告同预算下的准确率与延迟分布，并与单次生成、独立多采样和简单投票比较。

